\section{A Measurement Protocol for This Model}
\label{sec:protocol}

Everything in this section was found by getting it wrong first. Each item states the trap, the
failure mode it can produce, and the check that catches it. Historical pre-fix magnitudes are not
part of the submission artifact; the corrected protocol is the fact source for any future
measurement.

\subsection{Re-evaluate every non-seed node every round}

\paragraph{The trap.} Recording a node as ``settled'' the first time it activates. Under this model a
node's state is a function of the exposure it has received \emph{so far}, and exposure keeps growing
as its in-neighbours activate. A node that is positive in round $t$ can be over-exposed in round
$t+1$, and the model says it then becomes negative.

\paragraph{The error.} A counterexample with no Monte-Carlo noise: edges $s \to a$ with weight
$0.4$, $s \to b$ with weight $0.5$, $b \to a$ with weight $0.4$; seed $\{s\}$; windows
$a = [0.2, 0.6]$ and $b = [0.2, 0.9]$. Node $a$ first receives $\delta = 0.4$, inside its window, so
it turns positive. Node $b$ then activates and raises $a$'s exposure to $0.9 > 0.6$, so $a$ must
turn negative. Freezing $a$ gives spread $3$ and an empty negative set; the correct answer is spread
$2$ with $a$ negative.

\paragraph{Why it matters more than it looks.} On the audit graph the freezing error changes the
measured spread by far more than the policy differences that the paper is intended to compare.
Those pre-fix numerical measurements are omitted from the submission; the point retained here is
the mechanism, not a withdrawn benchmark.

\paragraph{The check.} The transition rule must be the literal one,
$\theta^\kappa_v \le \delta(v,t) \le \theta^\tau_v$, evaluated for every non-seed node in every round.
Two boundary cases catch most implementations: a node with $\delta = \theta^\tau = 1$ is
\emph{positive} (the comparison is inclusive), and a node whose $\delta$ exceeds $\theta^\tau$ after
having been positive must be moved out of the positive set.

\subsection{State the threshold law, because the sign of the effect depends on it}

\paragraph{The trap.} Sampling two independent uniforms and clamping the upper one to $1$, then
describing the result as ``linear thresholds''.

\paragraph{The error.} The model draws $(\theta^\kappa, \theta^\tau)$ uniformly from the simplex
$\{0 \le \theta^\kappa \le \theta^\tau \le 1\}$, i.e.\ as the order statistics of two independent
uniforms. The induced marginals are
\[
  F^\kappa(x) = 2x - x^2, \qquad F^\tau(x) = x^2,
\]
so $\mathbb{E}[\theta^\kappa] = 1/3$ and $\mathbb{E}[\theta^\tau] = 2/3$. Clamping
$\theta^\tau = 1$ keeps $F^\kappa = 2x - x^2$, but replacing the joint draw with
$\theta^\kappa \sim U[0,1]$ changes the lower marginal to $\mathbb{E}[\theta^\kappa] = 1/2$. The two
are different processes: on a one-hop instance with two candidates of weight $0.4$ into a single
target, the model's law gives $F_D(\{a\}) = 0.4800 > F_D(\{a,b\}) = 0.3200$ --- non-monotone ---
while uniform-threshold LT gives $0.4000 < 0.8000$, which is monotone. Every statement about
non-monotonicity, and therefore about the inapplicability of coverage sampling, holds for one law
and fails for the other.

\paragraph{The check.} Name the law in the configuration, and verify the realised marginals against
$2x - x^2$ and $x^2$ rather than trusting the sampler. Four laws must be distinguished, not two:
the simplex law; the same draw with $\theta^\tau$ clamped (over-exposure off, $\kappa$ marginal
preserved); genuine uniform-threshold LT; and the simplex law with $\theta^\tau$'s support raised,
which is an intervention knob rather than a degeneracy.

\subsection{Do not quote a rank correlation below $\mathrm{MC} = 300$ in the saturated regime}

\paragraph{The trap.} Measuring candidate marginals with a few dozen Monte-Carlo trials and
reporting rank correlations, or a share of candidates with a negative marginal.

\paragraph{The error.} In the regime this paper is about, the true marginals can be small compared
with the error of estimating them. Low-trial labels can therefore change sign and rank ordering as
the Monte-Carlo budget increases. The pre-fix numerical sweep is withdrawn; only the measurement
lesson is retained here.

\paragraph{The check.} Report the paired standard error of every label. A candidate is negative only
if its mean is more than two paired standard errors below zero, and a share of negative candidates
is only interpretable with that column present. Prefer paired spread comparisons, which cancel the
window-draw noise by construction and are unaffected by this trap.

\subsection{Never pool seed sizes in a state-dependence statistic}

\paragraph{The trap.} Measuring whether a candidate's marginal depends on the current seed set by
drawing contexts at several seed sizes --- typically $0, k, 2k, 3k$ --- and computing one
within/between variance ratio over all of them.

\paragraph{The error.} The unsaturated context can dominate the pooled statistic. A pooled ratio
therefore reports behaviour outside the regime of interest and can give the right qualitative
verdict for the wrong reason. We now compute the ratio inside one seed size and report it per size.

\paragraph{The check.} Compute the ratio inside one seed size and report it per size. A cell counts
as state dependent only if it clears the threshold at the seed size at which the rest of the
measurement was taken.

\subsection{Repeat the candidate pool, because one draw can flip the sign}

\paragraph{The trap.} Drawing one random pool of candidates and reporting the headroom a policy has
on it.

\paragraph{The error.} The pool's degree profile decides how much headroom is visible, and the
\emph{sign} of the policy gap can depend on the draw. A single-pool number therefore cannot support
a claim in either direction.

\paragraph{The check.} Draw the pool degree-stratified so its profile matches the graph's by
construction, repeat over several independent draws, and report the median with a min--max band.
A verdict should not be published from an unfinished sweep: if a run covers a subset of the
requested graphs, say so in the result file rather than reporting a partial aggregate.

\subsection{What these five checks cost}

They are cheap. Re-evaluating non-seed nodes adds one pass over the frontier per round. Naming the
threshold law is free. Storing a paired standard error alongside each label costs nothing, because
the paired differences are formed before averaging anyway. Computing the state-dependence ratio per
size rather than pooled is a loop boundary. Stratifying and repeating the pool multiplies one
measurement by the number of draws. The expensive check is the Monte-Carlo budget, which is a
factor of $25$ over a low-budget run.  The earlier low-budget tables are omitted from the
submission; the remaining experiment section records the protocol boundary and the pending
post-fix quality--cost study.
